\documentclass{article}
\usepackage{amsmath}
\usepackage{amsfonts}
\usepackage{amssymb}
\usepackage{graphicx}
\usepackage{hyperref}
\begin{document}

\noindent \textbf{Asymmetric Loss Function and $\lambda_{dd}$ Formalization.} 

To align the neural network's optimization objective with the real-world risk management goals of trend-following strategies, we introduce an asymmetric loss function. Standard Binary Cross-Entropy (BCE) loss treats False Positives (premature exits) and False Negatives (missed drawdown protections) symmetrically. However, in quantitative trading, the financial penalty of failing to capture a severe profit give-back is substantially higher than the opportunity cost of a slightly premature liquidation. 

Therefore, during the training phase, we employ a Weighted Binary Cross-Entropy loss governed by the drawdown penalty coefficient ($\lambda_{dd}$). Formally, for a given training sequence at time $t$, the loss function is defined as:

\begin{equation}
\mathcal{L} = - \frac{1}{N} \sum_{i=1}^{N} \left[ \lambda_{dd} \cdot y_i \log(p_i) + (1 - y_i) \log(1 - p_i) \right]
\end{equation}

where $y_i \in \{0, 1\}$ represents the binary supervisory label of early exit, and $p_i$ is the conditional probability predicted by the LSTM layer. 

The hyperparameter $\lambda_{dd} > 0$ acts as a financial risk-aversion multiplier. When $\lambda_{dd} > 1.0$, the gradient update heavily penalizes the model for failing to predict an early exit when a substantial trailing drawdown is imminent ($y_i = 1$). Economically, tweaking $\lambda_{dd}$ allows our framework to adapt to different market microstructures and accommodate traders with varying degrees of drawdown tolerance.

\vspace{0.5em}
\noindent \textbf{Empirical Discussion on Asset-Specific $\lambda_{dd}$ Variations.} 

As reported in Table 5, the optimal hyperparameter configurations derived via grid search exhibit intriguing cross-asset heterogeneity, particularly regarding the drawdown penalty coefficient $\lambda_{dd}$ \cite{source1}. The baseline asset, Bitcoin (BTC), yields an optimal $\lambda_{dd} = 1.0$ \cite{source1}. This implies that for a highly liquid and historically less volatile cryptocurrency, a symmetric or moderately balanced loss structure is sufficient for the LSTM model to map out stable exit boundaries. 

Conversely, for Ethereum (ETH) and Ripple (XRP), the optimal penalty climbs significantly to $\lambda_{dd} = 2.0$, while Solana (SOL) establishes an optimal threshold at $\lambda_{dd} = 1.5$ \cite{source1}. This phenomenological discrepancy can be rationalized by the inherent market microstructure and tail-risk characteristics of altcoins. High-beta assets such as SOL and XRP are notorious for their explosive late-stage trend accelerations followed by sharp, asymmetric mean-reversions—a primary driver of the severe profit give-back risk formalized in Section 1 \cite{source1}. For these assets, a standard loss function ($\lambda_{dd}=1.0$) would lean toward underestimating the swift destruction of paper profits. By allocating a higher loss weight ($\lambda_{dd} \ge 1.5$), the LSTM decision layer is forced to adopt a more conservative posture, prioritizing capital preservation and timely profit-locking, which empirically translates into superior Sharpe ratios and mitigated maximum drawdowns, as verified by our sensitivity analysis (see Figures 15–18 in Appendix B.3) \cite{source1}.

\begin{thebibliography}{99}
\bibitem{source1} Your reference here.
\end{thebibliography}

\end{document}